\let\negmedspace\undefined
\let\negthickspace\undefined
\documentclass[journal]{IEEEtran}
\usepackage[a5paper, margin=10mm, onecolumn]{geometry}
%\usepackage{lmodern} % Ensure lmodern is loaded for pdflatex
% \usepackage{tfrupee} % Include tfrupee package

\setlength{\headheight}{1cm} % Set the height of the header box
\setlength{\headsep}{0mm}     % Set the distance between the header box and the top of the text

\usepackage{gvv-book}
\usepackage{gvv}
\usepackage{cite}
\usepackage{amsmath,amssymb,amsfonts,amsthm}
\usepackage{algorithm}
\usepackage{algorithmic}
\usepackage{graphicx}
\usepackage{textcomp}
\usepackage{xcolor}
\usepackage{txfonts}
\usepackage{listings}
\usepackage{enumitem}
\usepackage{mathtools}
\usepackage{gensymb}
\usepackage{comment}
\usepackage[breaklinks=true]{hyperref}
\usepackage{tkz-euclide} 
\usepackage{listings}
% \usepackage{gvv}                                        
\def\inputGnumericTable{}                                 
\usepackage[latin1]{inputenc}                                
\usepackage{color}                                            
\usepackage{array}                                            
\usepackage{longtable}                
\usepackage{calc}                                             
\usepackage{multirow}                                         
\usepackage{hhline}                                           
\usepackage{ifthen}                             
\usepackage{caption}              
\usepackage{lscape}
\usepackage{subcaption}
\usepackage{circuitikz}
\captionsetup{compatibility=false}
% \captionsetup{compactibility=false}
% \usepackage{algpseudocode}
\begin{document}

\bibliographystyle{IEEEtran}
\vspace{3cm}

\title{Experiment 5} 
\author{S A Aravind Eswar and Eshan Sharma}
{\let\newpage\relax\maketitle}

\section{Aim}
Op-Amp Applications:
\begin{enumerate}
	\item Custom Weighted Summing and Difference Amplifier
	\item Op-Amp Integrator
	\item Precision Rectifier(Super Diode)
\end{enumerate}

\section{Materials and Apparatus Required}

\begin{enumerate}
	\item Op-Amp (e.g., LM386, LM58)
	\item Resistors(Carefully selected for weighting)
	\item Capacitors
	\item Diode(e.g., 1N4148)
	\item Function Generator
	\item DC Power Supply
	\item Oscilloscope
\end{enumerate}

\section{Theory}
	\subsection{Op-Amp Integrator}
	
	\subsection*{Circuit Description}
	The basic configuration of an op-amp integrator consists of:
	\begin{itemize}
		\item An operational amplifier.
		\item A resistor (\(R_{\text{in}}\)) connected to the inverting input.
		\item A capacitor (\(C\)) connected between the output and the inverting input.
		\item The non-inverting input is grounded.
	\end{itemize}
	
	The circuit diagram is shown below:
	\begin{figure}[!ht]
		\centering
		\resizebox{0.6\textwidth}{!}{%
			\begin{circuitikz}
				\tikzstyle{every node}=[font=\normalsize]
				\draw (6.25,14) to[R] (7.75,14);
				\draw (9,15.5) to[C] (11,15.5);
				\draw (10,13.5) node[op amp,scale=1, yscale=-1 ] (opamp2) {};
				\draw (opamp2.+) to[short] (8.5,14);
				\draw  (opamp2.-) to[short] (8.5,13);
				\draw (11.2,13.5) to[short](11.5,13.5);
				\draw (7.75,14) to[short] (8.5,14);
				\draw (8.5,14) to[short] (8.5,15.5);
				\draw (8.5,15.5) to[short] (9,15.5);
				\draw (11,15.5) to[short] (11.5,15.5);
				\draw (11.5,15.5) to[short] (11.5,13.5);
				\draw (11.5,13.5) to[short] (12.75,13.5);
				\draw (8.5,13) to[short] (8.5,12.5);
				\draw (5.5,14) to[short] (6.25,14);
				\draw (8.5,12.5) to (8.5,12.25) node[ground]{};
				\draw (5.5,13.25) to (5.5,12.75) node[ground]{};
				\draw (12.75,12.75) to (12.75,12.25) node[ground]{};
				\node at (8.5,14) [circ] {};
				\node at (11.5,13.5) [circ] {};
				\node at (5.5,14) [circ] {};
				\node at (5.5,13.25) [circ] {};
				\node at (12.75,13.5) [circ] {};
				\node at (12.75,12.75) [circ] {};
				\node [font=\normalsize] at (8.25,14.25) {X};
				\node [font=\normalsize] at (10.25,16) {C          };
				\node [font=\normalsize] at (6.75,14.5) {$R_{in}$};
				\node [font=\normalsize] at (13.25,13.25) {$-V_{out}$};
				\node [font=\normalsize] at (5,13.75) {$V_{in}$};
				\node [font=\normalsize] at (9.75,13.5) {A};
				\draw [->, >=Stealth] (5.5,14) -- (6.25,14);
				\draw [->, >=Stealth] (8.5,15.5) -- (9.25,15.5);
				\node [font=\normalsize] at (5.5,14.25) {$I_{in}$};
			\end{circuitikz}
		}%
		\label{fig:my_label}
	\end{figure}
	
	\subsection*{Operation of the Integrator Circuit}
	The operation of the integrator is based on the virtual ground principle at the inverting input of the op-amp. Assuming an ideal op-amp:
	\begin{enumerate}
		\item The input current through \(R_{\text{in}}\) is given by:
		\[
		I_{\text{in}} = \frac{V_{\text{in}}}{R_{\text{in}}}.
		\]
		\item Since no current flows into the op-amp's inverting terminal, this current flows through the capacitor, causing it to charge or discharge.
		\item The voltage across the capacitor (\(V_C\)) is related to the charge (\(Q\)) by:
		\[
		Q = C V_C.
		\]
		The current through the capacitor is:
		\[
		I_C = C \frac{\mathrm{d}V_C}{\mathrm{d}t}.
		\]
	\end{enumerate}
	
	Equating \(I_{\text{in}}\) and \(I_C\), we have:
	\[
	\frac{V_{\text{in}}}{R_{\text{in}}} = C \frac{\mathrm{d}V_{\text{out}}}{\mathrm{d}t}.
	\]
	
	Rearranging and integrating both sides with respect to time gives:
	\[
	V_{\text{out}} = -\frac{1}{R_{\text{in}}C} \int V_{\text{in}} \, \mathrm{d}t,
	\]
	where the negative sign indicates that the output is inverted.
	
	Thus, the output voltage (\(V_{\text{out}}\)) is proportional to the time integral of the input voltage (\(V_{\text{in}}\)).

	\subsection{Half and Full-Wave Rectifiers}
	
	\subsection*{Half-Wave Precision Rectifier}
	
	The half-wave precision rectifier allows only one half of the input AC signal (positive or negative) to pass through while blocking the other half. This is achieved using an op-amp and a diode configured such that the op-amp compensates for the diode's forward voltage drop.
	
	\subsubsection*{Circuit Description}
	\begin{enumerate}
		\item An operational amplifier (e.g., LM358).
		\item A diode (e.g., 1N4148) connected in feedback.
		\item Resistors to set the input and feedback paths.
	\end{enumerate}
	
	\begin{figure}[!ht]
		\centering
		\resizebox{0.6\textwidth}{!}{%
			\begin{circuitikz}
				\tikzstyle{every node}=[font=\normalsize]
				\draw (5.75,13.5) to[R] (7.25,13.5);
				\draw (9,15.5) to[R] (11.25,15.5);
				\draw (9,14.75) to[D] (11.25,14.75);
				\draw (11.25,13) to[D] (13,13);
				\draw (13.75,13) to[R] (13.75,11.25);
				\draw (7.25,13.5) to[short] (8.25,13.5);
				\draw (8.25,12.5) to[short] (8.25,11.75);
				\draw (8,13.5) to[short] (8,15.5);
				\draw (8,15.5) to[short] (9,15.5);
				\draw (8,14.75) to[short] (9,14.75);
				\draw (11.25,14.75) to[short] (11.25,13);
				\draw (11.25,15.5) to[short] (12.75,15.5);
				\draw (12.75,15.5) to[short] (12.75,13);
				\draw (13,13) to[short] (13.75,13);
				\draw (5,13.5) to[short] (6,13.5);
				\draw (5,12.5) to[short] (5.75,12.5);
				\draw (5.75,12.5) to[short] (5.75,11.75);
				\draw (5.75,11.75) to (5.75,11.5) node[ground]{};
				\draw (8.25,11.75) to (8.25,11.5) node[ground]{};
				\draw (13.75,11.25) to (13.75,11) node[ground]{};
				\node at (8,13.5) [circ] {};
				\node at (8,14.75) [circ] {};
				\node at (8,15.5) [circ] {};
				\node at (11.25,14.75) [circ] {};
				\node at (11.25,13) [circ] {};
				\node at (12.75,13) [circ] {};
				\node at (5,13.5) [circ] {};
				\node at (5,12.5) [circ] {};
				\draw (9.75,13) node[op amp,scale=1, yscale=-1 ] (opamp2) {};
				\draw (opamp2.+) to[short] (8.25,13.5);
				\draw  (opamp2.-) to[short] (8.25,12.5);
				\draw (10.95,13) to[short](11.25,13);
				\node [font=\normalsize] at (6.5,14) {9.1K};
				\node [font=\normalsize] at (10.25,16) {9.1K};
				\node [font=\normalsize] at (14.25,12.25) {1K};
				\node [font=\normalsize] at (4.5,13) {$V_{in}$};
				\node [font=\normalsize] at (14.25,13.25) {$V_{out}$};
			\end{circuitikz}
		}%
		
		\label{fig:my_label}
	\end{figure}
	
	\subsubsection*{Working Principle}
	When the input signal (\(V_{\text{in}}\)) is positive, the op-amp output drives the diode into conduction, and the output (\(V_{\text{out}}\)) follows \(V_{\text{in}}\). When \(V_{\text{in}}\) is negative, the diode becomes reverse-biased, blocking current flow, and \(V_{\text{out}}\) remains at zero.
	
	Mathematically:
	\[
	V_{\text{out}} = 
	\begin{cases} 
		V_{\text{in}}, & V_{\text{in}} > 0, \\
		0, & V_{\text{in}} \leq 0.
	\end{cases}
	\]
	
	This results in a rectified waveform where only the positive half-cycles are present at the output.
	
	\subsection*{Full-Wave Precision Rectifier}
	
	The full-wave precision rectifier allows both halves of the input AC signal to be rectified, producing a unipolar output. This is achieved by combining a half-wave rectifier with a summing amplifier.
	
	\subsubsection*{Circuit Description}
	\begin{enumerate}
		\item A half-wave precision rectifier circuit.
		\item A summing amplifier to combine signals from both positive and negative cycles.
		\item Additional diodes and resistors for proper signal routing.
	\end{enumerate}
	
	\begin{figure}[!ht]
		\centering
		\resizebox{0.7\textwidth}{!}{%
			\begin{circuitikz}
				\tikzstyle{every node}=[font=\normalsize]
				\draw (6.25,13.5) to[R] (7.75,13.5);
				\draw (9,15.5) to[R] (11.25,15.5);
				\draw (11.25,14.75) to[D] (9,14.75);
				\draw (13,13) to[D] (11.25,13);
				\draw (8,13.5) to[short] (8,15.5);
				\draw (8,15.5) to[short] (9,15.5);
				\draw (8,14.75) to[short] (9,14.75);
				\draw (11.25,14.75) to[short] (11.25,13);
				\draw (11.25,15.5) to[short] (12.75,15.5);
				\draw (12.75,15.5) to[short] (12.75,13);
				\draw (5,12.5) to[short] (5.75,12.5);
				\draw (5.75,12.5) to[short] (5.75,11.75);
				\draw (5.75,11.75) to (5.75,11.5) node[ground]{};
				\draw (8.25,12.5) to (8.25,12.25) node[ground]{};
				\node at (8,13.5) [circ] {};
				\node at (8,14.75) [circ] {};
				\node at (8,15.5) [circ] {};
				\node at (11.25,14.75) [circ] {};
				\node at (11.25,13) [circ] {};
				\node at (12.75,13) [circ] {};
				\node at (5,13.5) [circ] {};
				\node at (5,12.5) [circ] {};
				\draw (9.75,13) node[op amp,scale=1, yscale=-1 ] (opamp2) {};
				\draw (opamp2.+) to[short] (8.25,13.5);
				\draw  (opamp2.-) to[short] (8.25,12.5);
				\draw (10.95,13) to[short](11.25,13);
				\node [font=\normalsize] at (4.5,13) {$V_{in}$};
				\node [font=\normalsize] at (18.75,10) {$V_{out}$};
				\draw (9,11) to[R] (11,11);
				\draw (13,13) to[R] (14.5,13);
				\draw (7.75,13.5) to[short] (8.5,13.5);
				\draw (5,13.5) to[short] (6.25,13.5);
				\draw (16.25,10.5) node[op amp,scale=1, yscale=-1 ] (opamp2) {};
				\draw (opamp2.+) to[short] (14.75,11);
				\draw  (opamp2.-) to[short] (14.75,10);
				\draw (17.45,10.5) to[short](17.75,10.5);
				\draw (6.25,13.5) to[short] (6.25,11);
				\draw (6.25,11) to[short] (9,11);
				\draw (11,11) to[short] (14.75,11);
				\draw (14.5,13) to[short] (14.5,11);
				\draw (15,12.25) to[R] (17.25,12.25);
				\draw (15,12.25) to[short] (15,11);
				\draw (17.25,12.25) to[short] (17.25,10.5);
				\draw (14.75,10) to (14.75,9) node[ground]{};
				\draw (17.75,10.5) to[short] (18.75,10.5);
				\draw (17.75,9.5) to[short] (18.75,9.5);
				\draw (17.75,9.5) to (17.75,9) node[ground]{};
				\node at (14.5,11) [circ] {};
				\node at (15,11) [circ] {};
				\node at (17.25,10.5) [circ] {};
				\node at (18.75,10.5) [circ] {};
				\node at (18.75,9.5) [circ] {};
				\node [font=\normalsize] at (7,14) {6.8K};
				\node [font=\normalsize] at (10.25,16) {6.8K};
				\node [font=\normalsize] at (13.5,13.5) {3.4K};
				\node [font=\normalsize] at (16.25,12.75) {6.8K};
				\node [font=\normalsize] at (10,10.5) {6.8K};
			\end{circuitikz}
		}%
		\label{fig:my_label}
	\end{figure}
	
	\subsubsection*{Working Principle}
	\begin{enumerate}
	\item During the positive half-cycle of \(V_{\text{in}}\), the first op-amp operates as a half-wave rectifier, passing the positive signal to the summing amplifier.
	\item During the negative half-cycle of \(V_{\text{in}}\), another path inverts and shifts the signal such that it also appears as a positive contribution at the summing amplifier input.
	
	The summing amplifier combines these contributions to produce a continuous positive output for both halves of the input signal.
	
	Mathematically:
	For an input sinusoidal signal \(V_{\text{in}} = V_m \sin(\omega t)\),
	the output voltage is:
	\[
	V_{\text{out}} = |V_{\text{in}}| = V_m |\sin(\omega t)|.
	\]
	
	This results in a fully rectified waveform at the output.
	 \end{enumerate}

\section{Procedure}
	\subsection{Op-Amp Integrator}
	
\begin{enumerate}
	\item Set Up the Circuit:
\begin{itemize}
	\item Connect the circuit as per the integrator circuit diagram.
	\item Use a resistor (\(R_{\text{in}}\)) at the input and a capacitor (\(C_f\)) in the feedback loop of the op-amp.
	\item Ground the non-inverting terminal of the op-amp.
	\item Connect the power supply to the op-amp (\(12V\) and \(-12V\)).
\end{itemize}

\item Apply Input Signal:
\begin{itemize}
	\item Connect a function generator to provide a square wave input signal to the circuit.
\end{itemize}

\item Observe Output Waveform:
\begin{itemize}
	\item Connect the input and output signals to two channels of a CRO.
	\item Set the CRO to AC coupling mode.
	\item Observe and record both input and output waveforms simultaneously.
\end{itemize}

\item Analyze Output Waveform:
\begin{itemize}
	\item For a square wave input, verify that the output is a triangular waveform.
	\item Note down the amplitude and time period of both input and output signals.
\end{itemize}

\end{enumerate}

	\subsection{Precision Rectifiers}
	
	\begin{enumerate}
		\item Set Up the Circuit:
		\begin{itemize}
			\item Connect the circuit as per the precision rectifier circuit diagram.
			\item For the half-wave rectifier:
			\begin{itemize}
				\item Use an op-amp with a diode in the feedback loop and a resistor at the input.
				\item Ground the non-inverting terminal of the op-amp.
			\end{itemize}
			\item For the full-wave rectifier:
			\begin{itemize}
				\item Combine two op-amps: one for inverting and rectifying the negative half-cycle and another for summing both positive and inverted signals.
				\item Use appropriate diodes and resistors for signal routing.
			\end{itemize}
			\item Connect the power supply to the op-amp (\(+12V\) and \(-12V\)).
		\end{itemize}
		
		\item Apply Input Signal:
		\begin{itemize}
			\item Connect a function generator to provide a sine wave input signal to the circuit.
		\end{itemize}
		
		\item Observe Output Waveform:
		\begin{itemize}
			\item Connect the input and output signals to two channels of a CRO.
			\item Set the CRO to AC coupling mode.
			\item Observe and record both input and output waveforms simultaneously.
		\end{itemize}
		
		\item Analyze Output Waveform:
		\begin{itemize}
			\item For the half-wave rectifier:
			\begin{itemize}
				\item Verify that only one half (positive or negative) of the sine wave is present at the output while the other half is blocked.
				\item Note down the amplitude and time period of both input and output signals.
			\end{itemize}
			\item For the full-wave rectifier:
			\begin{itemize}
				\item Verify that both halves of the sine wave are rectified, resulting in a unipolar output waveform.
				\item Note down the amplitude and time period of both input and output signals.
			\end{itemize}
		\end{itemize}
	\end{enumerate}
\section{Observations}
	\subsection{Op-Amp Integrator}

\section{Precautions}

\begin{enumerate}
	\item Make sure not to accidentally discharge the capacitor
	\item Make sure the DC generator is connected to ground properly
	\item Use a push button for shorting the inductor and capacitor
	\item Repeat the experiment multiple times to be sure of the reading as it is a sensitive measurement
\end{enumerate}

\end{document}